\documentclass{article}
\usepackage{amsmath,amssymb,amsthm}
\newtheorem{lemma}{Lemma}
\newtheorem{theorem}{Theorem}
\begin{document}
\title{Actual-face overlap strengthens the full 683-edge degree gate}
\author{Near-extremal campaign, 2026-11-15}
\maketitle

\section{E01: Domain and exact imported interface}
Let $G\subseteq Q_8$ be C4-free with exactly 683 edges, and fix either
actual bipartition class $P$. The opposite class is unrestricted.
Let $n_d=|\{v\in P:d_G(v)=d\}|$, so $\sum n_d=128$ and
$\sum d n_d=683$. For each of the 112 coordinate $Q_6$ faces $F$, set
\[
 L_F=|\{v\in P\cap F:d_{G[F]}(v)\le2\}|,\quad
 S_F=\sum_{v\in P\cap F}\binom{d_{G[F]}(v)}3,\quad
 c_F=|\{v\in P\cap F:d_{G[F]}(v)=6\}|.
\]
The same graph and roots are used in every restriction.

We import the dimension-local bounds in the actual Cycle007 sources:
$2S_F+c_F\le316$ if $L_F=0$, $\le317$ if $L_F=1$, and
$\le321$ if $L_F=2$. Each permits arbitrary opposite degrees, including
zero. For $L_F=3$ the bound is 322: star injection gives $S_F\le160$,
full roots have $c_F\le4$, and
$S_F\equiv29+c_F\pmod3$ gives weighted caps $316,319,322,319,322$
for $c_F=0,\ldots,4$. For $L_F\ge4$, $S_F\le160,c_F\le4$ give 324.
Thus define
\[
 t(0)=0,\quad t(1)=1,\quad t(2)=5,\quad t(3)=6,
 \quad t(\ell)=8\quad(4\le\ell\le32),
\]
so every actual face satisfies $2S_F+c_F\le316+t(L_F)$.
These are proved local imports, not inequalities fitted to this campaign's
graphs. Their historical proof/replay limits remain explicit in the skeleton.

\section{E02: Equality retains the actual low-root overlap}
Set $b(d)=\binom d3-10d+40$ and $b_P=\sum_dn_db(d)$.
The specified-side star sum is $10\cdot683-5120+b_P$.
Every actual $Q_3$ star lies in ten $Q_6$ faces, and a global degree-$d$
root is full in $\binom d6$ such faces. Hence
\[
 Q:=\sum_F(2S_F+c_F)=20\sum_{v\in P}\binom{d_G(v)}3
                        +\sum_dn_d\binom d6.                 \tag{1}
\]
A root of degree $d$ is low in
\[
 \lambda(d)=\sum_{j=0}^2\binom dj\binom{8-d}{6-j}
 =(28,28,28,18,6,0,0,0,0)_d
\]
containing faces, so $\Lambda:=\sum_FL_F=\sum_dn_d\lambda(d)$.
Define the actual overlap loss
\[
 \Phi=\sum_F\left(\frac52 L_F-t(L_F)\right)\ge0.
\]
Its summand is zero exactly for $L_F\in\{0,2\}$, and is positive
otherwise (values $3/2,3/2$ at 1 and 3, and $5L_F/2-8\ge2$ at 4 or more).
Summing the local bounds gives
\[
 Q\le112\cdot316+\frac52\Lambda-\Phi.                       \tag{2}
\]
Using the two degree sums to eliminate $n_5,n_6$, (1)--(2) give
\[
 C(P)+\Phi\le1149,\qquad
 C(P)=735n_0+534n_1+333n_2+177n_3+66n_4+105n_7+345n_8.       \tag{3}
\]
For an explicit check, before subtracting $(5/2)\lambda(d)$ the
coefficients after elimination are
$(805,604,403,222,81,0,0,105,345)$, and
$138432-201\cdot683=1149$.

\begin{lemma}
If $C(P)=1149$, every actual $Q_6$ face has exactly zero or two low
specified-side roots. In particular every face containing a particular
low root must contain another distinct low root.
\end{lemma}
\begin{proof}
Equation (3) forces $\Phi=0$. All its summands are nonnegative and only
the values zero and two have zero summand. The final statement follows.
\end{proof}
No regularity, minimum degree or parity-square condition on $G$ was used.

\section{E03: Actual low-face sets and pair intersections}
For a global root $v\in P$ let $W(v)$ be the set of containing $Q_6$
faces on which it is low. Let $R(v)$ be its actual retained direction set.
A face is specified by an omitted coordinate pair $I$ and its two fixed
bits. It belongs to $W(v)$ exactly when those bits equal $v_I$ and
$|R(v)\setminus I|\le2$.
For degree at most two, $|W(v)|=28$; degree three gives 18, degree four
gives six, and degree at least five gives zero.

For distinct $u,v\in P$ the Hamming distance $r$ is 2,4,6 or 8.
A face can contain both only if $I$ avoids all differing coordinates,
so there are at most $\binom{8-r}{2}\le15$ common containing faces.
If $v$ has degree three, its omitted pair must meet $R(v)$. For $r=2$
there are six available coordinates, of which at most three belong to
$R(v)$, giving at most $\binom62-\binom32=12$ choices. For $r\ge4$
there are at most six choices in total. Thus
\[
 |W(u)\cap W(v)|\le12\quad\hbox{if }d(v)=3.                 \tag{4}
\]
If $d(v)=4$, the omitted pair must be contained in its four retained
directions, so the intersection is at most six. These bounds do not
assume that independently chosen stars extend to a full graph.

For degrees $i,j\in\{0,1,2,3,4\}$ the following uniform bounds therefore
apply to distinct same-parity roots:
\[
M_{ij}=\begin{cases}15&i,j\le2,\\12&\max(i,j)=3,\\6&\max(i,j)=4.
\end{cases}                                                    \tag{5}
\]
Put $\lambda_i=\lambda(i)$ and $k=\sum_{i=0}^4n_i$. A profile-only
lower bound $\Psi\le2\Phi$ is
\[
\Psi=\begin{cases}
0&k=0,\\
3\lambda_i&k=1\text{ with its low root of degree }i,\\
3(\lambda_i+\lambda_j-2M_{ij})&k=2,\\
3\max\left(0,\max_{i:n_i>0}\{\lambda_i-
                \sum_{j=0}^4 M_{ij}(n_j-\mathbf1_{i=j})\}\right)&k\ge3.
\end{cases}                                                    \tag{6}
\]
For one root all its low faces have $L_F=1$. For two roots the number
of such faces is $\lambda_i+\lambda_j-2|W(u)\cap W(v)|$ and all other
face counts are zero or two. For three or more roots, fix one of degree
$i$. At least $\lambda_i-\sum_jM_{ij}(n_j-\mathbf1_{i=j})$ of its low
faces contain no other low root, by the union bound on the actual
intersections. Each contributes $3/2$ to $\Phi$. Taking the maximum,
not the sum, proves (6). In particular every original graph satisfies
\[
                  2C(P)+\Psi\le2298.                         \tag{7}
\]

\section{E04: Five exact exclusions}
\begin{theorem}
The following specified-side profiles $(n_0,\ldots,n_8)$ cannot occur
in any C4-free 683-edge subgraph of $Q_8$:
\[
\begin{array}{c|rrrrrrrrr}
 A&0&0&0&1&2&86&31&8&0\\
 B&0&0&0&2&0&84&39&1&2\\
 C&0&0&1&0&6&73&44&4&0\\
 D&0&1&1&1&0&74&50&1&0\\
 E&1&0&0&1&2&73&50&1&0
\end{array}
\]
There is no restriction on the opposite-side degree profile.
\end{theorem}
\begin{proof}
Each row has 128 vertices, degree sum 683 and $C(P)=1149$; thus every
low-face count is zero or two by the lemma. Only the displayed roots of
degrees at most four can be low after restriction.

In A the degree-three root is low on 18 faces, but the two degree-four
roots supply only twelve other low incidences in total. They cannot give
a partner on every one of those 18 distinct faces.

In B both degree-three roots would have to be low on exactly the same
18 faces. Two distinct roots share at most fifteen containing faces,
a contradiction.

In D the degree-one root is low on all its 28 faces. Its intersections
with the degree-two and degree-three roots have sizes at most 15 and 12,
respectively, so at most 27 of those faces have a partner.
In E the degree-zero root similarly has 28 faces, while its intersections
with the degree-three and two degree-four roots have total at most
$12+6+6=24$. Both are impossible.

For C let $u$ be the unique degree-two root and $v_1,\ldots,v_6$ the
degree-four roots. Translate the entire graph by $u$, so $u=0$.
Write $R_j=R(v_j)$, each of size four. For any omitted pair $I$, summing
the even local low counts over its four fixed-bit faces shows
\[
 1+|\{j:I\subseteq R_j\}|\equiv0\pmod2.                    \tag{8}
\]
The one comes from $u$, which is always low; all other low incidences
are exactly those with $I\subseteq R_j$. Thus every coordinate pair
belongs to an odd number of the six four-sets.

Here is the complete consequence for those six sets, permitting repeated
sets initially. Let $r_a$ count sets containing coordinate $a$. Summing
(8) over the seven pairs containing $a$ gives $3r_a\equiv7\pmod2$,
so $r_a$ is odd. Covering all seven pairs also requires $3r_a\ge7$,
hence $r_a\ge3$. Since $\sum_a r_a=24$, every $r_a=3$.
Each coordinate therefore gives a three-subset of the six block indices.
For a coordinate $a$, its intersection sizes with the seven other
three-subsets are odd and total $3r_a=9$; exactly one is three and the
other six are one. Intersection three means identical three-subsets.
Consequently the eight coordinates split into four pairs $A,B,C,D$ of
identical incidence signatures. Each four-set $R_j$ consists of exactly
two such pairs. Each pair of different pairs occurs in precisely one
$R_j$, since their signatures intersect once. The six $R_j$ are exactly
the six unions of two of $A,B,C,D$.

For the unique block $R_j=A\cup B$, every omitted pair with one
coordinate in $A$ and one in $B$ is contained in no other $R_k$.
The face containing $u=0$ then needs $v_j$ as its only possible partner.
This forces $(v_j)_a=(v_j)_b=0$ for every such pair, hence $v_j$ is zero
on all four coordinates of $R_j$. The same holds for every block.
Now omit the two coordinates of $A$ and fix them to zero. This face
contains $u$ and the three degree-four roots belonging to $A\cup B$,
$A\cup C$ and $A\cup D$, all low there. It has at least four low roots,
contradicting the required zero-or-two count. This excludes C.
\end{proof}

\section{E05: Complete finite profile cover and exact filter}
Every original graph has nonnegative integer $n_0,\ldots,n_8$ with
the two degree sums, $C(P)\le1149$, and $b_P\le82$ (nonnegative global
star holes). Consider precisely this necessary integer relaxation,
without asserting that any of its points is a graph.
Enumerate $n_0,n_1,n_2,n_3,n_4,n_7,n_8$ using their positive coefficients
in $C(P)$. The remaining two values are uniquely
\[
 n_6=683-\sum_{d\ne5,6}dn_d-5\left(128-\sum_{d\ne5,6}n_d\right),
 \quad n_5=128-\sum_{d\ne5,6}n_d-n_6.
\]
Keeping their nonnegative values and $b_P\le82$ is complete for the
declared relaxation. The exact enumeration gives 814 profiles.
Equation (7) rejects sixteen; the separate six-block argument for C
rejects one more not rejected by (7). Thus 797 profiles remain in this
specified necessary cover. The complete input and output rows are saved
in PROFILES.json, and an independently structured finite checker must
verify these counts before computational acceptance. The written proof
establishes why every discarded row is impossible and why every original
graph maps into the enumerated inputs.

\section{E06: Main-target connection and limits}
Every arbitrary 683-edge C4-free Q8 graph has both parity profiles
subject to (3), including all degrees zero through eight. The theorem
therefore removes these five profiles on either side of the full domain;
it does not assume them for all candidates or discard their complement.
The previous odd-small-defect gate still permits restricting the original
683 problem to at least nineteen even squares, at its recorded import
and same-family assurance. No square-parity hypothesis was needed here.
Neither these exclusions nor an unsuccessful solver search proves the
unrestricted upper bound 682. No literature priority is claimed.
The count797 describes only the explicit one-sided integer relaxation,
not a new graph classification or a claim that all previous independent
restrictions have been intersected. The actual-face inequality (3) can
be stricter than its profile-only consequence (7).
\end{document}
